% Datos de los autores
\newcommand{\getTitle}{Arquitectura de una plataforma educativa basada en videojuegos adaptativos con control docente para el aprendizaje de la programación}
\newcommand{\getAuthor}{Johny Antonio Pedraza Romero}

\title{\getTitle}
\author{\getAuthor}
\newcommand{\grade}{...}
\newcommand{\authorEmail}{jpedraza@ceis.cujae.edu.cu} 
\newcommand{\supervisorName}{DrC. Raisa Socorro Llanes}
\newcommand{\supervisorEmail}{rsocorro@ceis.cujae.edu.cu} 

% Datos de la Universidad
\newcommand{\university}{Universidad Tecnológica de La Habana "José Antonio Echeverría" CUJAE}
\newcommand{\faculty}{Facultad de Informática}
\newcommand{\carrer}{Ingeniería Informática}

% Datos del documento
\newcommand{\chapterOne}{Marco teorico y estado del arte}
\newcommand{\chapterSecond}{Diseño de la solución}
\newcommand{\chapterThree}{Validación de la solución}

\renewcommand{\appendixname}{Anexo}

% Funciones propias
\newcommand{\testcase}[3]{
	\begingroup
	\begin{table}[H]
		\centering
		\caption{#1} % Parámetro 1: El título de la tabla
		\label{#2}   % Parámetro 2: El identificador para referenciar (\ref{})
		\renewcommand{\arraystretch}{1.4} 
		
		\begin{tabularx}{\textwidth}{ | >{\bfseries}p{4.5cm} | X | }
			\hline
			#3 % Parámetro 3: Todo el contenido interno de la tabla
			\hline
		\end{tabularx}
	\end{table}
	\endgroup
}


% Función para generar tablas de Descripción de Casos de Uso
% Parámetro 1: Nombre del Caso de Uso (para el título)
% Parámetro 2: Label para referenciar (\ref{})
% Parámetro 3: Contenido de la tabla
\newcommand{\usecase}[3]{
	\begingroup
	\begin{table}[H]
		\centering
		\caption{#1} 
		\label{#2}   
		\renewcommand{\arraystretch}{1.3} 
		
		% m{3.5cm} -> Ancho fijo para los nombres de los campos
		% X -> El resto del ancho de la página para la descripción
		% >{\bfseries} -> Pone en negrita automático la primera columna
		\begin{tabularx}{\textwidth}{ | >{\bfseries}m{3.5cm} | X | }
			\hline
			#3 
			\hline
		\end{tabularx}
	\end{table}
	\endgroup
}

\NewDocumentCommand{\tablaCasoPrueba}{m m m m +m}{%
	% ---- Bloque de encabezado ----
	\noindent
	\begin{tblr}{
			width   = \linewidth,
			colspec = {Q[l,font=\bfseries,wd=3.8cm] X[l]},
			hlines, vlines,
			rowsep  = 3pt,
		}
		\SetCell[c=2]{l} \textbf{Caso de Prueba} & \\
		Nombre                  & #1 \\
		Objetivo de la Prueba   & #2 \\
		Descripción de la prueba & #3 \\
		Condiciones de entrada  & #4 \\
		\SetCell[c=2]{l} \textbf{Combinaciones de valores} & \\
	\end{tblr}
	
	% ---- Bloque de combinaciones ----
	\noindent
	\begin{tblr}{
			width   = \linewidth,
			colspec = {Q[c,wd=0.6cm] X[2,l] X[1,l] X[3,l] X[2,l] X[2,l]},
			hlines, vlines,
			rowsep  = 3pt,
		}
		\textbf{CP} &
		\textbf{Escenarios} &
		\textbf{Variables} &
		\textbf{Valores de entrada} &
		\textbf{Resultado esperado} &
		\textbf{Resultado real} \\
		#5
	\end{tblr}
}
	